% ================= TÓM TẮT =================
\newpage
\thispagestyle{plain}

\begin{center}
    {\bfseries\Large TÓM TẮT}
\end{center}

\vspace{1cm}

Phát hiện biển báo giao thông là một khối chức năng bắt buộc của hệ thống hỗ trợ lái tiên tiến và
xe tự hành. Điểm khó của bài toán không nằm ở việc đạt một chỉ số mAP cao trên tập kiểm tra, mà ở
chỗ mô hình phải chạy được trên phần cứng nhúng gắn trên xe, nơi bộ nhớ và ngân sách năng lượng
đều bị giới hạn nghiêm ngặt. Điều này dẫn tới nhu cầu \emph{nén mô hình}: làm cho mô hình nhỏ hơn
và nhanh hơn mà không đánh mất năng lực phát hiện. Hai kỹ thuật nén phổ biến nhất là chưng cất tri
thức và lượng tử hoá sau huấn luyện. Đề tài đặt câu hỏi hai kỹ thuật này thực sự mang lại gì trong
bối cảnh một bộ dữ liệu biển báo quy mô vừa, và quan trọng hơn, cái giá phải trả nằm ở đâu.

Đề tài sử dụng bộ dữ liệu \textit{pkdarabi/cardetection} trên Kaggle, gồm 4.969 ảnh độ phân giải
416$\times$416 với 6.012 hộp bao thuộc 15 lớp là đèn tín hiệu và biển giới hạn tốc độ. Trước khi
huấn luyện, một phân tích khám phá đầy đủ được thực hiện: phân bố lớp, thống kê kích thước hộp bao,
bản đồ nhiệt vị trí vật thể, kiểm tra chất lượng nhãn và kiểm toán rò rỉ dữ liệu giữa các tập. Năm
cấu hình được khảo sát trong cùng một lượt chạy: mô hình thầy YOLO26s, mô hình học trò YOLO26n
huấn luyện thông thường làm đối chứng, mô hình học trò được chưng cất từ thầy, cùng hai bản xuất
sang ONNX ở độ chính xác FP32 và INT8.

Điểm then chốt về mặt phương pháp là toàn bộ năm cấu hình được đánh giá bằng \textbf{pycocotools}
trên cùng một tệp chú thích COCO, nhờ đó ngoài mAP tổng còn thu được $AP$ tách riêng cho từng nhóm
kích thước vật thể. Ngoài ra, năm ngưỡng chấp nhận được khai báo \emph{trước} khi chạy thực nghiệm,
buộc người thực hiện cam kết tiêu chí đánh giá trước khi nhìn thấy số liệu.

Kết quả cho thấy mô hình học trò được chưng cất đạt mAP@0,5 bằng 0,9014, \emph{thấp hơn} mô hình
đối chứng huấn luyện thông thường (0,9348) 3,34 điểm phần trăm --- một kết quả âm đi ngược giả
thuyết ban đầu. Lượng tử hoá INT8 giảm dung lượng tệp 70,3\% và chỉ làm mAP@0,5 giảm 0,0038. Tuy
nhiên chỉ số tách theo kích thước vật thể cho một bức tranh hoàn toàn khác: cùng phép lượng tử hoá
đó làm $AP$ trên nhóm vật thể nhỏ giảm 0,0247, tức \textbf{gấp 6,5 lần} mức suy giảm của mAP tổng.
Trong năm ngưỡng chấp nhận, ba đạt và hai không đạt, cả hai ngưỡng không đạt đều thuộc về giả
thuyết chưng cất.

Luận điểm này được kiểm chứng theo hai hướng độc lập. Hướng thứ nhất loại bỏ khả năng kết quả là
tạo tác của dữ liệu bẩn: khi đo lại chỉ trên 573 ảnh kiểm tra không bị rò rỉ, tỉ số giữa mức sụt
$AP$ nhỏ và mức sụt mAP tổng ở phép chưng cất \emph{tăng} từ 1,85 lên 2,66 lần. Hướng thứ hai dùng
ba đoạn video đường phố ngoài phân bố huấn luyện. Vì video không có nhãn thật nên phép kiểm chứng
không đo độ chính xác, mà kiểm tra một dự đoán có thể sai: nếu điểm yếu thật nằm ở vật thể nhỏ thì
việc tăng độ phân giải đầu vào phải làm lộ thêm các phát hiện nhỏ một cách chọn lọc. Kết quả đúng
như dự đoán --- mức tăng số phát hiện giảm đơn điệu theo kích thước, lần lượt 13,0 --- 4,3 --- 2,3
lần ở ba nhóm nhỏ dần, và \emph{giảm} một nửa ở nhóm lớn nhất.

Đóng góp của đề tài không nằm ở một con số mAP cao mà ở hai phát hiện mang tính phương pháp luận.
Thứ nhất, mAP tổng che giấu gần như hoàn toàn chi phí thật của các phép nén mô hình; chỉ khi tách
$AP$ theo nhóm kích thước mới thấy được rằng toàn bộ suy giảm dồn vào vật thể nhỏ --- đúng nhóm
quan trọng nhất với bài toán biển báo, nơi biển ở xa quyết định thời gian phản ứng của xe. Thứ hai,
phân tích thành phần bộ dữ liệu chỉ ra 38,68\% số ảnh là ảnh cắt cận cảnh kiểu GTSRB khiến phân bố
kích thước bị lưỡng cực, và kiểm toán rò rỉ phát hiện 10,2\% tập kiểm tra có bản sao trong tập huấn
luyện --- trong đó 71,3\% số hộp bao bị rò rỉ thuộc đúng nhóm vật thể nhỏ. Hạn chế lớn nhất là nhánh
chưng cất chạy với kích thước lô nhỏ hơn nhánh đối chứng, một yếu tố gây nhiễu chưa được kiểm soát.

\vspace{0.5cm}
\noindent\textbf{Từ khoá:} phát hiện biển báo giao thông, chưng cất tri thức, lượng tử hoá,
nén mô hình, vật thể nhỏ, kết quả âm, rò rỉ dữ liệu.
